\section{Conclusion}
\label{sec:conclusion}

Color recovery for multiplexed QR codes does not need offline training. Each photograph is
fitted on its own function patterns: a closed-form inverse of the leak between edge-adjacent
modules, with one fixed fraction $\beta=0.06$, followed by regularized discriminant analysis
with shrink $\lambda=0.5$. These values were chosen on each dataset separately; they are the
best setting on the qrism dataset and statistically tied with the best on CUHK-CQRC
($p\HiqGridBestP$). At these values the method decodes \HiqCodeChosen\% of localized symbols
on CUHK-CQRC and \QrCodeChosen\% on the qrism dataset, and every baseline decodes fewer on
both. Given the same module samples on HiQ's test split, it decodes \PortOursCode\% of symbols
against \PortCmiCode\% for the best of HiQ's trained classifiers, at \TimeChosenNs\,ns per
module including the fit.

Color codes have been proposed for file transfer between phones without a
network~\cite{toh2016colorqr}, for validation data on small product labels such as medicine
packs that must be checked off-line~\cite{melgar2019hd2dc}, for encrypted payloads carried
with their signatures~\cite{wahsheh2026rgb}, and for hidden traceability and
anti-counterfeiting layers beside a layer that standard readers
decode~\cite{huang2022mc2d,huang2023halftone}. In these uses the printer, display and phone
are often not known in advance. Our method assumes nothing about them, only that the
function patterns are intact: it is fitted to each photograph without training data,
recovers the colors of a version-40 symbol in \TimeChosenVforty\,ms, and was evaluated on
both printed and displayed symbols. The decode rates above are for single photographs; a
phone scanner sees many frames per second and could retry on later ones, but we have not
measured how often this helps.

The fixed $\beta$ is the main compromise: its best value shifts with module size, capture
distance and lighting, and choosing it per photograph in a way that tracks decodability rather than
bit error rate remains open.
The Rust implementation is available at \url{https://github.com/mohnishsalian0/qrism}, and
the qrism dataset is released with its source renders.
